\scanpage{266}
\mspage{78}
% archive sheet number: 266
une façon ``canonique'', ``naturelle'', ``voulue par Dieu'', de la ``\unsure{remplumer}'' en un diagramme substantiel.

Les flèches les plus évidentes qui se \deleted{\unsure{présentent}} sont \deleted{\unsure{les}} relatives à un foncteur

\[
f:I\unsure{\text{---}|\longrightarrow}J
\]

de $\mathrm{Dia}$, et aux deux morphismes d'adjonction pour le couple $(f^*,f_*)$, savoir

% slanted
\margin{Bien sûr, problème dual pour $(f_!,f^*)$}

\[
\begin{aligned}
G&\longrightarrow f_*f^*G&&\text{dans }\mathbb D(J)\\
f^*f_*F&\unsure{\text{---}|\longrightarrow}F&&\text{dans }\mathbb D(I).
\end{aligned}
\]

Je vais essayer de remplumer ces deux flèches, et vais généraliser en \deleted{\illegible{2}}\inserted{envisageant la} situation ci-contre, où

\diagram{\begin{tikzcd}[ampersand replacement=\&,row sep=large,column sep=large]
I\arrow[d,"f"']\&F\\
J\&G\arrow[u,"u"']
\end{tikzcd}}

\[
\operatorname{Hom}\unsure{{}^{f}|f|{}_{f}}(G,F)\overset{\mathrm{def}}{=}\operatorname{Hom}(f^*G,F)
\overset{\mathrm{adj}}{\simeq}\operatorname{Hom}(G,\inserted{f_*(F)})
\]

Donc le $=$ $u$ s'interprète par \uline{deux} flèches squelettiques

\[
\left\{\begin{aligned}
f^*(G)&\longrightarrow F&&\text{dans }\mathbb D(I)\\
G&\longrightarrow f_*(F)&&\text{dans }\mathbb D(J).
\end{aligned}\right.
\tag*{(1)}
\]

Les remplumer, c'est les interpréter (\unsure{se|si} faire se peut) en termes de deux objets (si possible canoniques) de $\mathbb D(I\times\mathbb{\Delta}_1)$ et $\mathbb D(J\times\deleted{\unsure{\mathbb D}}\mathbb{\Delta}_1)$ respectivement.

Je vais me guider sur le cas particulier d'un dérivateur trivial

\[
\mathbb D=\mathbb D_{\mathcal M},
\]

où $\mathcal M$ est une catégorie de modèles ``nue'' ($\Sigma=\varnothing$).
